\newpage
\section{有关三角函数与解三角形必背的几个公式}
\begingroup
\setlength{\abovedisplayskip}{4pt}
\setlength{\belowdisplayskip}{4pt}
\setlength{\abovedisplayshortskip}{2pt}
\setlength{\belowdisplayshortskip}{2pt}
这些公式要求会默写、会推导，并能够在题目中灵活应用。

\paragraph{1°诱导公式} \mbox{}
\begin{equation*}
\sin \left( \alpha + 2 k \pi \right) = \sin \alpha , \cos \left( \alpha + 2 k \pi \right) = \cos \alpha , \tan \left( \alpha + k \pi \right) = \tan \alpha , k \in \mathbb{Z} .
\end{equation*}
\begin{equation*}
\sin \left( \alpha + \pi \right) = - \sin \alpha , \cos \left( \alpha + \pi \right) = - \cos \alpha , \tan \left( \alpha + \pi \right) = \tan \alpha
\end{equation*}
\begin{equation*}
\sin \left( - \alpha \right) = - \sin \alpha , \cos \left( - \alpha \right) = \cos \alpha , \tan \left( - \alpha \right) = - \tan \alpha
\end{equation*}
\begin{equation*}
\sin \left( \pi - \alpha \right) = \sin \alpha , \cos \left( \pi - \alpha \right) = - \cos \alpha , \tan \left( \pi - \alpha \right) = - \tan \alpha
\end{equation*}
\begin{equation*}
\sin \left( \frac{\pi}{2} - \alpha \right) = \cos \alpha , \cos \left( \frac{\pi}{2} - \alpha \right) = \sin \alpha , \sin \left( \frac{\pi}{2} + \alpha \right) = \cos \alpha , \cos \left( \frac{\pi}{2} + \alpha \right) = - \sin \alpha
\end{equation*}
\paragraph{2°和差角公式、二倍角公式} \mbox{}
\begin{equation*}
\sin \left( \alpha + \beta \right) = \sin \alpha \cos \beta + \cos \alpha \sin \beta , \quad \cos \left( \alpha + \beta \right) = \cos \alpha \cos \beta - \sin \alpha \sin \beta
\end{equation*}
\begin{equation*}
\sin \left( \alpha - \beta \right) = \sin \alpha \cos \beta - \cos \alpha \sin \beta , \quad \cos \left( \alpha - \beta \right) = \cos \alpha \cos \beta + \sin \alpha \sin \beta
\end{equation*}
\begin{equation*}
\tan \left( \alpha + \beta \right) = \frac{\tan \alpha + \tan \beta}{1 - \tan \alpha \tan \beta} , \tan \left( \alpha - \beta \right) = \frac{\tan \alpha - \tan \beta}{1 + \tan \alpha \tan \beta}
\end{equation*}

请直接写出二倍角公式。

\paragraph{3°和差化积/积化和差公式及其他二级公式} \mbox{}
\begin{equation*}
\sin \alpha + \sin \beta = 2 \sin \frac{\alpha + \beta}{2} \cos \frac{\alpha - \beta}{2} , \quad \sin \alpha - \sin \beta = 2 \sin \frac{\alpha - \beta}{2} \cos \frac{\alpha + \beta}{2}
\end{equation*}
\begin{equation*}
\cos \alpha + \cos \beta = 2 \cos \frac{\alpha + \beta}{2} \cos \frac{\alpha - \beta}{2} , \quad \cos \alpha - \cos \beta = -2 \sin \frac{\alpha + \beta}{2} \sin \frac{\alpha - \beta}{2}
\end{equation*}
\begin{equation*}
\tan \alpha + \tan \beta = \frac{\sin \left( \alpha + \beta \right)}{\cos \alpha \cos \beta} , \quad \tan \alpha - \tan \beta = \frac{\sin \left( \alpha - \beta \right)}{\cos \alpha \cos \beta}
\end{equation*}
\begin{equation*}
\cos^{2} \beta - \cos^{2} \alpha = \sin^{2} \alpha - \sin^{2} \beta = \sin \left( \alpha + \beta \right) \sin \left( \alpha - \beta \right)
\end{equation*}
\begin{equation*}
\cos^{2} \alpha - \sin^{2} \beta = \cos^{2} \beta - \sin^{2} \alpha = \cos \left( \alpha + \beta \right) \cos \left( \alpha - \beta \right)
\end{equation*}
\begin{equation*}
\sin \alpha \cos \beta = \frac{1}{2} \left[ \sin \left( \alpha + \beta \right) + \sin \left( \alpha - \beta \right) \right] , \quad \cos \alpha \sin \beta = \frac{1}{2} \left[ \sin \left( \alpha + \beta \right) - \sin \left( \alpha - \beta \right) \right]
\end{equation*}
\begin{equation*}
\cos \alpha \cos \beta = \frac{1}{2} \left[ \cos \left( \alpha + \beta \right) + \cos \left( \alpha - \beta \right) \right] , \quad \sin \alpha \sin \beta = \frac{1}{2} \left[ \cos \left( \alpha - \beta \right) - \cos \left( \alpha + \beta \right) \right]
\end{equation*}

\paragraph{4°正/余弦定理} \mbox{}

已知三角形\(A B C\)，角\(A , B , C\)的对边分别为\(a , b , c\)，则下式恒成立：
\begin{equation*}
\frac{a}{\sin A} = \frac{b}{\sin B} = \frac{c}{\sin C}
\end{equation*}
\begin{equation*}
a^{2} = b^{2} + c^{2} - 2 b c \cos A , \quad b^{2} = a^{2} + c^{2} - 2 a c \cos B , \quad c^{2} = a^{2} + b^{2} - 2 a b \cos C
\end{equation*}
\endgroup
